% IgboAI: a living resource/infrastructure paper.
% Versioning: \paperversion below tracks the repo release tag; the PDF is
% built by CI (.github/workflows/paper-build.yml) and attached to releases.
\documentclass[11pt]{article}

\usepackage[margin=1in]{geometry}
\usepackage[T1]{fontenc}
\usepackage[utf8]{inputenc}
\usepackage{booktabs}
\usepackage{graphicx}
\usepackage{hyperref}
\usepackage{xcolor}
\hypersetup{colorlinks=true, linkcolor=blue!60!black, urlcolor=blue!60!black, citecolor=blue!60!black}

\newcommand{\paperversion}{v0.1.0-draft}
\newcommand{\todo}[1]{\textcolor{red!70!black}{[TODO: #1]}}

\title{IgboAI: Self-Updating Research Infrastructure for a Low-Resource,
Tonal Language\\[0.4em]\large Version \paperversion}
\author{Ignatius Ezeani\thanks{\todo{affiliation, email, ORCID; add
co-authors as they join}}}
\date{\today}

\begin{document}
\maketitle

\begin{abstract}
\todo{150--200 words once Sections 1--3 stabilise.}
We describe IgboAI, an open, continuously self-updating research infrastructure for Igbo language technology: automated literature tracking,
licensing-aware corpus ingestion, and (in progress) benchmarking, built as LLM-agent workflows gated by human review. We report the design, the
practical challenges of low-resource data collection, and analyses of the collected material.
\end{abstract}

\section{Introduction}
\todo{Aims of the IgboAI project; why Igbo (speaker population, tonality, orthographic history, resource scarcity); why infrastructure-first rather
than model-first; relationship to the African NLP community.}

\section{Related Work}
\todo{African/Igbo NLP resources and surveys; Masakhane-line participatory research; automated literature tracking and living reviews; data governance
and data statements for low-resource corpora. Seed from \texttt{related-work/entries.jsonl}.}

\section{System Architecture}
\label{sec:architecture}
\todo{The two-layer principle: deterministic fetch/compute, LLM curation and documentation, human judgment at every merge. Describe the literature
tracker (nightly; arXiv, OpenAlex, HF) and the corpus pipeline (weekly; per-source storage policies). Include the layer-ownership model: policy
files humans write and robots read; data cards robots draft and humans approve; raw data only the deterministic layer touches. Diagram.}

\subsection{Safeguards and Economics}
\todo{Bounded work per run (item caps, turn caps, spend caps); audit trails (discard lists, seen-ledgers); branch protection and least-privilege
credentials; cost figures per run and per month.}

\section{Challenges of Low-Resource Ingestion}
\label{sec:challenges}
\todo{The honest chapter. Draft from \texttt{paper/notes/drudgery.md}: the failure ladder to a first green run; licensing triage per source class
(why news text is manifest-only; the JW300 precedent for religious text; unlicensed GitHub repositories); quality problems in Igbo Wikipedia
(alphabetical-backfill skew, markup leakage, MT-suspect orthography, archaic register) and the flagging policy adopted; the review-time
economics of human-in-the-loop curation.}

\section{Corpus Analysis}
\label{sec:corpus}
\todo{Auto-generated statistics will be included from \texttt{paper/stats.tex} once the derivation tooling lands.}
\IfFileExists{stats.tex}{\input{stats.tex}}{\todo{stats.tex not yet generated; tables of corpus size, source composition, flag rates, and
literature-tracker throughput will appear here.}}

\section{Benchmarks and Models}
\todo{Phase 4--5 of the roadmap: benchmark harness design, leaderboard, baseline results as they arrive. Placeholder until then.}

\section{Limitations and Ethical Considerations}
\todo{What automation cannot judge (register, dialect, cultural context); licensing residual risk; representation and community authorship; the
environmental/cost footprint of agentic workflows.}

\section{Conclusion}
\todo{One paragraph; forward pointer to the roadmap.}

\paragraph{Acknowledgements.}\todo{community, funders, infrastructure.}

\paragraph{Availability.} All code, data cards, and this paper are versioned at \url{https://github.com/IgboNLP-Research/IgboAI}; releases carry DOIs via Zenodo. \todo{insert concept DOI after first release.}

% \bibliographystyle{plainnat}
% \bibliography{references}
\end{document}
